\chapter{Discussion, Limitations and Future Work}\label{cha:discussion}
\label{cha:conclusion}

\section{Summary and Diagnosis}
\label{sec:disc-diagnosis}\label{sec:concl-summary}

DIPhysTwin adapts PhysTwin to multi-view RGB video: camera poses, depth and ground alignment are
inferred, then handed to its reconstruction and warp-model stages. The new component is a residual
depth smoother, supervised by Video Depth Anything (VDA) as a motion oracle and trained without
ground-truth depth \parencite{phystwin, vda}.

The pipeline is still 2D-first: segmentation and CoTracker run on RGB, and depth lifts their output
into 3D \parencite{cotracker}. Depth errors therefore reach the tracks, the shape prior, the
reconstruction and the warp model at once. No tested depth model was good at both jobs. DUSt3R,
MASt3R and DA3 give geometry but flicker, while VDA is smooth but monocular and geometrically weak
\parencite{dust3r, mast3r, da3, vda}. The smoother takes geometry from the first group and motion
from the second, at the cost of some blur (Chapter~\ref{cha:evaluation}).

Two alternatives were considered and dropped. Replacing the weak parts of the PhysTwin pipeline
with a multi-pose streaming depth model and multi-view tracking would require a new model
\parencite{dust3r, mast3r, cotracker}. Rewriting around 4D Gaussian splatting would remove explicit
poses but move segmentation and tracking into 3D \parencite{fourdgs, gaussiansplatting}.
\label{sec:disc-alternatives}

\section{Limitations}
\label{sec:disc-limitations}

The depth stage is not truly streaming. Duplicating fixed views as a dynamic sequence helps only
over short windows; larger chunks make the Umeyama alignment singular \parencite{umeyama1991}. The
three cameras also differ in intrinsics and distortion, use rolling shutters and automatic exposure
and focus, and were never time-synchronised; software only aligned their start times. The later rig
used measured intrinsics and extrinsics only as a reference, and may still be slightly fisheye. These conditions explain
the observed symptoms: flicker, spatial drift within a clip, and failures on fast motion and unusual
rotations. \label{sec:disc-streaming}\label{sec:disc-data}

The environment is static. The floor collider fixed the over-stiff behaviour for supported
objects, but a rope has little ground support and still needs unrealistic stiffness. The smoother
blurs depth, and the current evaluation cannot say whether it removes flicker without also removing
fine geometry. \label{sec:disc-environment}\label{sec:disc-smoothing}

The depth comparison covers one sequence, the remaining evidence is qualitative, and the loss
weights are untuned. The findings on object size, speed and occlusion are observed trends from a small number of sessions.
\label{sec:disc-validity}

\section{Future Work}
\label{sec:concl-future}

The depth stage needs true multi-view streaming. A model that accepts several camera poses would
replace the chunk-overlap workaround and its unstable alignment \parencite{umeyama1991}. The same
model should be both spatially consistent and temporally smooth, and the shape prior needs to handle
non-watertight meshes \parencite{arap}. \label{sec:concl-streaming}\label{sec:concl-depth}

The capture should be repeated with synchronisation cues, locked exposure and focus, and measured
intrinsics, allowing for fisheye distortion.
Tracking needs to survive occlusion and partially visible objects; MVTracker and direct 3D tracking
are the next candidates \parencite{mvtracker, spatrackerv2}. Environment interaction
should move from static planes to moving rigid bodies. \label{sec:concl-data}\label{sec:concl-tracking}\label{sec:concl-environment}

The smoother needs loss-term and span ablations, multi-view evaluation, and comparison against
reference geometry where any exists. After that the pipeline should be cleaned up and released.
\label{sec:concl-code}

\section{Conclusion}
\label{sec:concl-closing}

DIPhysTwin turns multi-view RGB video into an interactive digital twin, and the smoother reduces
temporal flicker without ground-truth depth. Results are strongest for large, slow-moving
objects under controlled capture. Depth coherence, occlusion, calibration and tracking remain the
main limitations.
